\section{Monolith 与 Blast Radius：代码边界不等于故障边界}

Asif 并不把 microservice 当作规模化前提。课程更关注 \term{blast radius}（爆炸半径）：一次代码、配置或依赖故障能影响多少用户、功能、region 和时间。Monolith 可以共享代码和部署流程，同时通过 traffic cell、queue、database、feature flag 与权限边界隔离高风险路径。

\subsection{Critical path 与 experimental path 分开}

本节回答“一个仓库怎样仍然拥有小故障域”。读图时不要数 service 数量，而要检查 auth、editor request、index job、new model 与 shared schema 的运行边界。若 experimental code 可以耗尽 critical database connection 或污染共享 queue，即使它被拆成独立 service，blast radius 仍然很大。

\lecturefigure{03-blast-radius.png}{Monolith 可以通过部署 cell、依赖和流量边界缩小 blast radius。}{本地字幕 05:20--07:50；概念重绘。}

\begin{knowledgebox}{读图：四个隔离旋钮}
\textbf{Traffic} 用 canary、tenant/region cell；\textbf{resource} 用独立 quota、queue 和 pool；\textbf{state} 用 versioned schema 与 rebuildable derived data；\textbf{change} 用 feature flag、rollback 与 deploy policy。Microservice 只是实现这些隔离的一种方式。
\end{knowledgebox}

\begin{importantbox}{Blast radius 的诊断式}
可粗略写为
\[
B=A\times D\times S,
\]
其中 $A$ 是 affected requests/users，$D$ 是持续时间，$S$ 是状态损坏或恢复复杂度。架构应优先降低不可逆的 $S$，再通过 canary 和 cell 降低 $A$，通过 detection/rollback 降低 $D$。
\end{importantbox}

\begin{warningbox}{“保持代码简单”不是少写抽象}
简单性应体现在 invariant 可解释、dependency 少、failure mode 可测试和 recovery path 明确。把复杂性藏进 global state、隐式 retry 或共享 database，代码行数可能减少，系统复杂度反而上升。
\end{warningbox}

\teachervoice{Asif 的选择并非“永远不要拆服务”，而是优先让少量工程师理解完整 critical path。只有当组织、故障域或资源形状真的需要独立演进时，再为分布式边界支付成本。}

\subsection{本章小结}

规模化首先是隔离问题，而不是 service 数量问题。好的 monolith 与好的 distributed system 都需要显式 blast radius、资源预算、版本契约和恢复策略。

